\section{Linear estimates on separated boundary faces}
\label{sec:linear-estimates}

The local estimates below record the precise classical boundary inputs used in the scalar continuation arguments.

\begin{lemma}[Scalar estimates on separated boundary faces]
\label{lem:linear-separated-faces}
Let $U\Subset V$ be nested interior or boundary patches with uniformly
controlled smooth geometry. Let $g$ be uniformly positive definite with
uniformly controlled $C^2$ coordinate norms. Suppose that $u$ is smooth,
\[
 Pu=-\Delta_g u+B\cdot du+cu=f,
 \qquad \|B\|_\infty+\|c\|_\infty\le K,
\]
and that the physical boundary in $V$, if present, carries either
$u=0$ or $\partial_{\nu_g}u=0$. For $1<p<\infty$,
\[
 \|u\|_{W^{2,p}(U)}
 \le C\bigl(\|f\|_{L^p(V)}+\|u\|_{L^p(V)}\bigr).
\]
The constant depends on the patch separation, ellipticity, principal
coefficient modulus, and $K$, but not on a continuity modulus of $B$.
If the coefficients and right side have uniform $C^{0,\gamma}$ bounds,
the corresponding one-sided $C^{2,\gamma}$ estimates hold; Dirichlet
data have norm $C^{2,\gamma}$ and Neumann data norm $C^{1,\gamma}$.

On a compact connected smooth domain, let the Dirichlet and Neumann
faces be disjoint unions of entire boundary components. With H\"older
coefficients, $c\ge0$, and either $c\ge c_*>0$ or a nonempty Dirichlet
face, $P$ is an isomorphism from the $C^{2,\gamma}$ space with homogeneous
boundary conditions to $C^{0,\gamma}$.
The Dirichlet Schauder a priori estimate also holds for a general
uniformly elliptic principal matrix $A\in C^{0,\gamma}$, with
H\"older lower-order coefficients. If $A\in C^{1,\gamma}$, the same
isomorphism conclusion holds with all-Dirichlet boundary conditions.
\end{lemma}

\begin{proof}
In metric boundary-normal coordinates,
$g=dt^2+h_{ab}(y,t)dy^a dy^b$ and the Neumann condition is $u_t=0$.
Reflect $u,h^{ab}$, tangential drift, $c$, and $f$ evenly across $t=0$,
and normal drift oddly. The principal coefficients remain Lipschitz
and uniformly elliptic. Normal drift may jump but is bounded. The
reflected first derivatives agree in trace, so the extended function
has weak second derivatives in $L^p$, with no measure on the plane,
and satisfies the extended equation almost everywhere. For zero
Dirichlet data use odd reflection of $u,f$ and the same coefficient
reflections; the same trace argument applies.

The interior estimate used on the reflected ball follows from the
constant-coefficient estimate of Agmon--Douglis--Nirenberg,
Theorem~14.1$'$, p.~700~\cite{agmondouglisnirenberg1959}.
For a cutoff function $v$ supported where
$C_0\operatorname{osc}A\le1/2$, freeze $A$ at $x_0$ to obtain
\[
 \|D^2v\|_p
 \le C_0\|A(x_0):D^2v\|_p
 \le C_0\|A:D^2v\|_p+\tfrac12\|D^2v\|_p.
\]
The cutoff commutators and bounded drift involve only $u,Du$.
Apply $\|Du\|_p\le\varepsilon\|D^2u\|_p+C_\varepsilon\|u\|_p$
on nested patches, absorb, and cover by finitely many small patches.
This proves the stated estimate, first for smooth functions and then
for strong $W^{2,p}$ solutions by approximation.

The one-sided Schauder estimate is the scalar specialization of
Agmon--Douglis--Nirenberg, Theorem~7.3, pp.~667--669; the oblique
estimate is also stated in Lieberman, Lemma~1,
p.~755~\cite{lieberman1982}. At a flattened face the normal
characteristic root $\tau_+$ has positive imaginary part. The
Dirichlet symbol is $1$; an oblique symbol has value
$\beta_T\cdot\xi+\beta_n\tau_+$, which cannot vanish since
$\beta_n\ne0$. Thus the required complementing condition holds.
The boundary faces never meet, so a finite cover uses one boundary
type in each boundary chart. Cutoffs constant in the normal variable
near a Neumann face preserve its homogeneous data. No estimate at an
artificial half-ball rim is used.

For solvability, add a sufficiently large constant $\kappa$.
The weak form for $P+\kappa$ on $H^1$ functions with zero Dirichlet
trace is coercive, since
\[
 \left|\int(B\cdot du)u\right|
 \le\tfrac12\|du\|_2^2+C\|B\|_\infty^2\|u\|_2^2.
\]
It therefore has a unique weak solution. For smooth data, tangential
difference quotients, followed by recovery of the second normal
derivative from the equation, give boundary regularity; iteration
gives smoothness. Approximation, the maximum principle for
$P+\kappa$, and the Schauder estimate give $C^{2,\gamma}$ solvability
for H\"older data. The original map $P$ is a compact perturbation of
$P+\kappa$ and hence has index zero. Its kernel vanishes by the
maximum principle and Hopf Boundary Lemma; constants are excluded by
the positive potential or by the nonempty Dirichlet face. Thus $P$
is invertible. Smooth inhomogeneous data are handled by a fixed lift.
The analogous all-Dirichlet argument for a general $C^{1,\gamma}$
positive principal tensor follows by rewriting it in divergence form,
putting its coefficient derivatives in the drift, and repeating the
shift and compact-perturbation argument.
\end{proof}

\paragraph{Application to the preparation equation.}
After $u=v/\eta_R$, the gradient term is
$D\sqrt{\zeta^2+|du|^2}$, bounded by $D(\zeta+|du|)$.
Here $\zeta=1$ in the stress preparation and $\zeta=w_2$ in
Equation~\eqref{eq:ah-conformal-corrector}; the positive smooth
regularizing weight has controlled norms on each fixed patch.
The preceding $W^{2,p}$ estimate and first-derivative interpolation
therefore bound $u$ using its height and the controlled source.
Take $p>3/(1-\gamma)$ for a permitted $0<\gamma<1$.
Sobolev embedding gives a H\"older gradient, so the original smooth
gradient nonlinearity is H\"older and the one-sided Schauder estimate
applies. One does not apply Schauder to the reflected equation with
its possibly discontinuous normal drift. On physical end patches,
multiply the estimate by the weight at the patch center; comparability
on the larger patch gives the weighted estimate without altering
homogeneous boundary data. All constants are uniform when the local
geometry and coefficient bounds are uniform.
